\section{Further research questions and an order of attack}
\label{research:section}

The following questions ask for exact identities, finite reductions,
or normalized analytic representations. They are not asserted
conjectures unless an implication is explicitly proved above.
The requested deliverable for each question is stated to separate
mathematical progress from additional numerical evidence.

\subsection{Higher digamma powers and harmonic Laurent coefficients}

\paragraph{1. The fifth power with all finite corrections.}
Repeat the quartic contour and finite summation argument for
$\FP\int_0^1\psi(x)^5\,\dd x$. The principal parts are
algorithmically finite harmonic polynomials, but their integration
requires a new exact summation-by-parts reduction. Determine
which linear Laurent coefficients of strict depth-four and
lower-depth zeta functions enter, and give every ordinary zeta
derivative and Stieltjes correction. A proof should include the
constant term of each finite asymptotic sum, not only matching
principal parts or a high-precision decimal.

\paragraph{2. Canonical polynomials that remove lower-depth data.}
The polynomial $T^4+4\gamma T^3$ removes $\eta$ from the quartic
bridge. Is there a systematically constructed monic polynomial
$P_k(T)$, with coefficients in a stated ordinary-constant algebra,
for which $\FP\int_0^1P_k(\psi(x))\,\dd x$ retains only one
new highest-depth harmonic coefficient? The Gamma and inverse-Gamma
coefficient polynomials are natural candidates for organizing the
recursion. A useful theorem would give the recurrence and prove
the cancellation for arbitrary $k$; the single quartic cancellation
does not establish such a theorem.

\paragraph{3. The arithmetic content of the depth-three coordinate.}
The normalized elementary kernel \eqref{q4:theta-kernel} determines
$\theta$ exactly. Seek an independent identity placing it in a
smaller algebra of zeta derivatives, ordinary Stieltjes constants,
or Gamma/Barnes integrals. The target should state which coordinates
are eliminated. Merely replacing the strict Laurent definition
by another name for the same integral does not reduce the
arithmetic content. Conversely, a failed numerical relation search
does not prove irreducibility.

\paragraph{4. A common-Hurwitz-shift quartic bridge.}
For $a>0$, define a consistent common shift in every denominator
of the strict depth-three series and derive the counterpart of
\eqref{q4:Didentity}. Determine the associated identity for
products of $\psi(a+x)$ and its derivatives, including the two
ordinary endpoint values on $[a,a+1]$. Then analyze the singular
limit $a\downarrow0$ with a specified coordinate finite part.
Replacing $\gamma_m$ by $\gamma_m(a)$ in the unshifted formula
without deriving these endpoint terms is not valid evidence.

\subsection{Harmonic-tail arithmetic beyond two factors}

\paragraph{5. The remaining harmonic factors in R1.}
For the products
\[
A_p(n)A_q(n)H_n^a
\prod_{j\ge1}\bigl(H_n^{(2j+1)}\bigr)^{b_j},
\]
with finitely many nonzero $b_j$, determine finite formulas for
the regular centered values and a moment-independent coordinate
space for each fixed polynomial. The source parity theorem predicts
local regularity, but it does not evaluate the global finite
constant. Finite summation by parts increases the depth of the
harmonic sums. One must specify the necessary multiple-zeta
coordinates and prove any claimed depth reduction.

\paragraph{6. Four even tails and the first genuine depth increase.}
Carry out the analogous calculation for
$A_pA_qA_rA_s$, first for $A_2^4$ and $A_2^3A_4$.
An even number of even tails preserves the centered even-power
asymptotics. The two-tail theorem suggests a finite Bernoulli
reduction at every moment order, but higher-depth convergent sums
can survive. Determine a fixed spanning set and an analytic
generating function. A parity argument alone does not prove that
all such higher-depth values reduce to ordinary zeta products.

\paragraph{7. First spectral derivatives at the regular lattice.}
Find exact reductions of $D'_{p,q}(-2r)$ or combinations of these
derivatives. Differentiating the subtracted continuation produces
logarithmically weighted global remainder sums as well as
derivatives of the explicitly removed Hurwitz terms. Seek
identities of these kernels that cancel the global terms, or
identify a minimal additional harmonic-Stieltjes coordinate
family. The finite value formula cannot simply be differentiated
with respect to its discrete integer label $r$.

\paragraph{8. Finite generation by initial moments.}
For fixed even $p,q$, Corollary~\ref{tail:fixed-space} gives at
most $\max(p,q)$ specified arithmetic coordinates. Determine
whether a uniform finite prefix of the moment sequence generates
the same rational span, and exhibit a rational inverse matrix when
it does. The $(2,4)$ case is proved by its first four moments.
An optimal bound concerns the rank of a specified coefficient
matrix. Equality with a dimension of actual constants would require
an additional arithmetic independence result.

\paragraph{9. Rational shifts and colored polylogarithms.}
Replace the inner shift $n+1$ and the center $n+1/2$ by compatible
rational shifts. Classify when the removed polynomial sums vanish,
and determine the surviving finite corrections when they do not.
The natural targets are Hurwitz zeta values at rational arguments
and colored Euler sums, together with polylogarithms at roots of
unity. A paired formulation at shifts $a$ and $1-a$ may retain
more symmetry than an isolated noncentral shift.

\subsection{Resolvent products, spectral derivatives, and primitives}

\paragraph{10. Several distinct nonreal resolvents.}
For
\[
\prod_{j=1}^d R_{z_j}(x)^{\lambda_j},
\]
derive an explicit joint Fourier or beta transform and complete
the hyperplanes where the total endpoint exponent is integral.
When all $z_j$ lie in the same half-plane the disk factors have
a common Fourier orientation; an Appell or Lauricella function
is a plausible ordinary-moment generator. Parameters in opposite
half-planes introduce two-sided convolutions and additional
arithmetic structure. Each claim needs a branch prescription and
the full endpoint amplitude before any specialization.

\paragraph{11. Products of continuous polygamma orders.}
The transform in this article is linear in the Hurwitz factor.
For products of two independently varying Hurwitz orders, derive
the convolution and identify exactly when ordinary polylogarithm
order derivatives suffice and when ordered or Tornheim coordinates
remain. Mixed order derivatives should be extracted from a joint
completed germ. Multiplying two separately completed scalar
identities does not account for the product's local singularities.

\paragraph{12. Arithmetic at rational complex-power exponents.}
At $q=0$, every Bernoulli moment is a finite polynomial in the
generalized harmonic values $H_\lambda^{(j)}$ and the fixed
prefactor. At rational $\lambda$ these values have classical
Hurwitz-zeta descriptions. Determine comparable finite reductions
for useful nonzero algebraic $q$ and rational $\lambda$, especially
for derivatives in $\lambda$ and the polylogarithm order.
The beta--polylogarithm formula supplies a convergent starting
integral; it does not by itself prove a finite constant evaluation.

\paragraph{13. Normalized primitives beyond the balanced combinations.}
Combine the explicit $z$-ladder with argument integration and
Hurwitz order differentiation to seek individual Barnes-type
representations. State the additive normalization at a base point
and derive its transport under $z$ and argument shifts. The
existing difference primitives already fix their constants; a
new result should identify a smaller named function class or an
additional functional equation, rather than introduce an arbitrary
integration constant.

\subsection{The Gaussian bridge and exact certification}

\paragraph{14. A functional identity for the frozen Gaussian residual.}
Express the actual $S_6$ or revised $S_8$ residual as a coefficient
of a proved joint generating identity, with all changes of
variables and branches explicit. The Cayley-ideal obstruction
specifies information that its named relation family cannot
supply. A successful new relation should have a computable
nonzero action on the separating functional, or an analytic
specialization demonstrably outside that family. The article's
resolvent transform is a source of candidate relations, not an
already established bridge.

\paragraph{15. Complete and certify any enlarged relation search.}
If the larger modular calculation is resumed, first preserve the
frozen target, relation enumeration, and checkpoint hashes. Finish
the declared relation system and recover an exact rational
certificate or exact separating functional. A modular rank
observation alone is not a rational proof of the target identity,
and an interrupted elimination does not show nonmembership in
the generated span. The input system and the terminal conclusion
must be independently replayable.

\paragraph{16. Formalize reusable analytic and algebraic components.}
The most reusable formal units are the two-point
Mittag--Leffler subtraction, the finite Faulhaber constant-term
calculation, the beta coefficient identity
\eqref{power:beta-coeff}, and the moving-monomial completion lemma
inside Theorem~\ref{power:completion}. Formalizing these would
make subsequent all-index identities considerably easier to
verify. An interval evaluation of $\theta$ or $Q_4$ can be added
separately by enclosing the same analytic tail formulas; it
would strengthen the numerical artifact, not replace the exact
proof.

\subsection{Practical sequence}

The most direct next calculations are the fifth-power reduction,
the initial-moment coefficient matrices, and the first additional
harmonic factors in the two-tail family. Each has a finite exact
starting point supplied here. Four-tail products and joint
resolvent powers are natural next structural extensions. The
Gaussian reduction should proceed only with a concrete candidate
functional bridge or a completed exact relation computation;
more precision for the same unproved scalar fit does not remove
the known missing step.

